% ZPD veidne pēc "Vadlīnijas skolēnu zinātniskās pētniecības darbu izstrādei un vērtēšanai" (2026).
% Kompilē: latexmk -pdf template.tex   (vai: make template no mapes zpd_darbs/)
% Sarkanās piezīmes [ ... ] paskaidro, kas katrā vietā jāraksta. Gatavam darbam: \documentclass[nolabels]{zpd}
\documentclass{zpd}
\graphicspath{{./}{../figures/}}

\skola{[Izglītības iestādes nosaukums]}
\nosaukums{[Darba nosaukums, ne vairāk kā 5 vārdi]}
\nozaru{[nozaru grupas / nozares] nozarē}
\autors[Darba autors]{[Vārds Uzvārds]}
\vaditajs[Darba vadītājs]{[Vārds Uzvārds]}
\konsultants[Darba konsultants]{[Vārds Uzvārds]}
\pilseta{[Pilsēta]}
\gads{[Gads]}

\begin{document}

\titullapa

% ------------------------------------------------------------------ anotācijas: abas vienā lapā
\anotacijulapa
\begin{anotacija}{Anotācija}
\zpdlabel{Autors un darba nosaukums}[Vārds Uzvārds]. [Darba nosaukums].

\zpdlabel{Darba saturs}[Darba mērķis, kā veikts pētījums, galvenie rezultāti un secinājumi. Kopā ne vairāk par pusi lappuses.]

\atslegvardi{[ne vairāk kā 5 atslēgvārdi, kas neatkārto nosaukumu]}
\end{anotacija}

\begin{anotacija}{Abstract}
\zpdlabel{Anotācijas precīzs tulkojums svešvalodā}[Author]. [Title].

[Translation of the annotation.]

\atslegvardi[Keywords]{[up to five keywords]}
\end{anotacija}

% ------------------------------------------------------------------ saturs
\saturs

% ------------------------------------------------------------------ ievads: no šejienes redzami lappušu numuri
\pamatteksts
\nenumuretanodala{Ievads}
\zpdlabel{Temata izvēles pamatojums}[Kāpēc izvēlēts šis temats.]

\zpdlabel{Pētījuma problēma un aktualitāte}[Kāda problēma tiek risināta un kāpēc tā ir aktuāla.]

\zpdlabel{Darba mērķis}[Viens mērķis, īsi un precīzi.]

\zpdlabel{Darba uzdevumi (3--5)}
\begin{enumerate}
  \item {}[Uzdevums.]
  \item {}[Uzdevums.]
  \item {}[Uzdevums.]
\end{enumerate}

\zpdlabel{Hipotēze vai pētījuma jautājums}[Hipotēze apgalvojuma formā vai atvērts jautājums (nesākas ar "vai").]

\zpdlabel{Pētījuma metodes}[Izmantotās metodes.]

% ------------------------------------------------------------------ satura nodaļas
\section{[Literatūras apskats]}
\zpdnote{Teorētiskā daļa: vismaz divu zinātnieku darbi šajā jomā, salīdzināti un kritiski izvērtēti. Ne vairāk kā 1/3 no darba teksta. Nodaļai nevar būt tikai viena apakšnodaļa.}
\subsection{[Apakšnodaļa]}
[Teksts ar atsauci \atsauce{Uzvārds, 2024, 5}.]
\subsection{[Apakšnodaļa]}
[Teksts.]

\zpdlabel{Nodaļas kopsavilkums}[Viena vai divas rindkopas un pāreja uz nākamo nodaļu.]

\section{[Metodes]}
\zpdnote{Datu vākšanas un analīzes metodes, to izvēles pamatojums, rīki, procedūra un ierobežojumi.}
[Teksts.]

\begin{equation}
  [formula] \label{eq:piemers}
\end{equation}

\section{[Rezultāti un analīze]}
\zpdnote{Visu iegūto datu apkopojums tabulās un attēlos, to analīze un salīdzinājums ar citu autoru datiem. Uz katru attēlu un tabulu tekstā ir vismaz viena norāde; vienā lappusē ne vairāk kā divi.}
[Teksts, sk. \ref{tab:piemers}~tabulu un \ref{fig:piemers}~attēlu.]

\begin{table}[H]
  \caption{[Tabulas nosaukums, saprotams bez teksta]}\label{tab:piemers}
  \centering
  \begin{tabular}{lrr}
    \toprule
    {}[Kolonna] & [Kolonna] & [Kolonna] \\
    \midrule
    {}[a] & [1] & [2] \\
    \bottomrule
  \end{tabular}
\end{table}

\begin{figure}[H]
  \centering
  \fbox{\parbox[c][3cm][c]{8cm}{\centering [attēls]}}
  \caption{[Attēla nosaukums, saprotams bez teksta]}\label{fig:piemers}
\end{figure}

% ------------------------------------------------------------------ secinājumi
\nenumuretanodala{Secinājumi}
\zpdnote{Tēžu veidā, numurēti, ne vairāk kā uzdevumu; atbild uz pētījuma jautājumu vai apstiprina/noraida hipotēzi.}
\begin{secinajumi}
  \item {}[Secinājums par 1. uzdevumu.]
  \item {}[Secinājums par 2. uzdevumu.]
  \item {}[Secinājums par 3. uzdevumu.]
\end{secinajumi}

% ------------------------------------------------------------------ avoti (5. pielikums vadlīnijās)
\begin{avoti}
  \avots Uzvārds, I. (Gads). \textit{Grāmatas nosaukums}. Vieta: Izdevējs. 100 lpp.
  \avots Uzvārds, I., Uzvārds, I. (Gads). Raksta nosaukums. \textit{Žurnāla nosaukums}, Vol. 1, pp. 1--10.
  \avots Uzvārds, I. (Gads). \textit{Nosaukums} [tiešsaiste]. Vietne [skatīts 2026. gada 1. novembrī]. Pieejams: \url{https://example.org}
  \avots Izstrādātājs. (Gads, datums). \textit{Rīka nosaukums, versija un sarunas nosaukums} [tiešsaiste]. [skatīts ...]. Pieejams: \url{https://example.org}
\end{avoti}

% ------------------------------------------------------------------ pielikumi (ne vairāk kā 1/3 no darba apjoma)
\pielikums{[Pielikuma nosaukums]}
[Saturs, uz kuru ir norāde tekstā.]

\end{document}
